\begin{proposition}[6.11.2]
\label{IV.2.s6.11.2-fr}
Sea $X$ un preesquema localmente noetheriano y localmente inmersible en
un esquema regular \textup{(5.11.1)}; sea $\mathcal F$ un $\mathcal O_X$-Módulo
coherente. Entonces:
\begin{enumerate}
\item[\textup{(i)}] \textup{(M. Auslander)} La función
$x\leadsto\operatorname{coprof}(\mathcal F_x)$ es semicontinua
superiormente en $X$ \textup{(es decir, para todo entero $n$, el conjunto
$U_{C_n}(\mathcal F)$ de los $x\in X$ tales que
$\operatorname{coprof}(\mathcal F_x)\leq n$ es} abierto).
\item[\textup{(ii)}] Para todo entero $n$, el conjunto $U_{S_n}(\mathcal F)$ de los
$x\in X$ tales que $\mathcal F$ posee la propiedad $(S_n)$ en el punto $x$
es abierto.
\end{enumerate}
\end{proposition}

\begin{proof}
\begin{enumerate}
\item[(i)] Puesto que la cuestión es local en $X$, es posible, en virtud de
la hipótesis, limitarse al caso en que $X$ es un subesquema cerrado de un esquema
afín regular $Y$. Si $j:X\to Y$ es la inyección canónica y
$\mathcal G=j_*(\mathcal F)$, se sabe que, para todo $x\in X$, se tiene
$\dim(\mathcal F_x)=\dim(\mathcal G_x)$ y
$\operatorname{prof}(\mathcal F_x)=\operatorname{prof}(\mathcal G_x)$,
por tanto, $\operatorname{coprof}(\mathcal F_x)=
\operatorname{coprof}(\mathcal G_x)$; como
$\operatorname{coprof}(\mathcal G_x)=0$ para $y\in Y\setmin X$, la cuestión se reduce a
probar la propiedad para $\mathcal G$. Es decir, es posible limitarse al
caso en que $X$ es un esquema afín \emph{regular}. Se sabe entonces (\textbf 0, 17.3.4) que
se tiene
\[
\operatorname{prof}(\mathcal F_x)=\dim(\mathcal O_{X,x})\setmin
\operatorname{dim.proj}(\mathcal F_x).
\]
Por otra parte, si $S$ es el único subpreesquema cerrado de $X$ cuyo
espacio subyacente es $\operatorname{Supp}(\mathcal F)$, se tiene
$\dim(\mathcal F_x)=\operatorname{codim}(\overline{\{x\}},S)$ (5.1.12.1) y, en
virtud de (5.1.9),
\[
\dim(\mathcal F_x)=\dim(\mathcal O_{X,x})\setmin\operatorname{codim}_x(S,X)
\]
pues $\mathcal O_{X,x}$ es un anillo regular y \emph{a fortiori}
biequidimensional (\textbf 0, 16.5.12). Se puede, por tanto, escribir
\begin{equation}
\operatorname{coprof}(\mathcal F_x)=
\operatorname{dim.proj}(\mathcal F_x)\setmin\operatorname{codim}_x(S,X)
\tag{6.11.2.1}\label{IV.2.s6.11.2.1-fr}
\end{equation}
y la proposición se sigue de (6.11.1) y de (\textbf 0, 14.2.6).

\item[(ii)] Como los $U_{C_n}(\mathcal F)$ son abiertos para todo $n$ por
(i), los $Z_n=X\setmin U_{C_n}(\mathcal F)$ son cerrados en $X$; además, es
claro que $Z_{n+1}\subset Z_n$ y, como $\dim(\mathcal F_x)$ es finita
para todo $x\in X$ y
$\operatorname{coprof}(\mathcal F_x)\leq\dim(\mathcal F_x)$,
la intersección de los $Z_n$ es \emph{vacía}; como es posible limitarse al caso en que $X$
es afín y, por tanto, cuasicompacto, se puede suponer que existe un $m$
\oldpage[IV-2]{160}
tal que $Z_m=\varnothing$. Ahora bien, se sigue de (5.7.4) que la relación
$x\in U_{S_k}(\mathcal F)$ equivale al conjunto de relaciones
\begin{equation}
\operatorname{codim}_x(Z_n,S)>n+k
\tag{6.11.2.2}\label{IV.2.s6.11.2.2-fr}
\end{equation}
para todo $n\geq0$; pero, para $n\geq m$, esta relación se verifica
automáticamente, de modo que basta considerar las relaciones (6.11.2.2)
para $0\leq n<m$. Ahora bien (\textbf 0, 14.2.6), el conjunto $V_{n,k}$ de los $x$ que verifican
(6.11.2.2) es abierto, y $U_{S_k}(\mathcal F)$, intersección de los $V_{n,k}$
para $0\leq n<m$, es también abierto. Queda demostrado.
\end{enumerate}
\end{proof}

Recordemos que la hipótesis de que $X$ es localmente inmersible en un esquema
regular siempre se cumple cuando $X$ es un preesquema \emph{localmente de
tipo finito sobre un cuerpo} $k$ (5.8.3).

\begin{corollary}[6.11.3]
\label{IV.2.s6.11.3-fr}
Bajo las hipótesis de \textup{(6.11.2)}, el conjunto $\operatorname{CM}(\mathcal F)$
de los puntos $x\in X$ tales que $\mathcal F_x$ es un módulo de
Cohen-Macaulay es abierto en $X$.
\end{corollary}

\begin{proof}
En efecto, es el conjunto $U_{C_0}(\mathcal F)$.
\end{proof}

\begin{remarks}[6.11.4]
\label{IV.2.s6.11.4-fr}
\begin{enumerate}
\item[(i)] El razonamiento de (6.11.2, (ii)) prueba que (sin hipótesis sobre
$X$), si $U_{C_n}(\mathcal F)$ es abierto para \emph{todo} entero $n$, entonces
$U_{S_n}(\mathcal F)$ es abierto para todo entero $n$.
\item[(ii)] Se tiene
$\operatorname{CM}(\mathcal F)=\bigcap_nU_{S_n}(\mathcal F)$. Si todo punto
$x\in X$ admite una vecindad abierta $V$ de dimensión \emph{finita}, la
sucesión de intersecciones $V\cap U_{S_n}(\mathcal F)$ es \emph{estacionaria},
pues existe un $m$ tal que $\dim(\mathcal F_x)\leq m$ para todo
$x\in V$; por consiguiente, si los $U_{S_n}(\mathcal F)$ son abiertos para \emph{todo}
entero $n$, $\operatorname{CM}(\mathcal F)$ es abierto en $X$.
\item[(iii)] Se escribirá $U_{S_n}(X)$, $U_{C_n}(X)$,
$\operatorname{CM}(X)$ en lugar de $U_{S_n}(\mathcal O_X)$,
$U_{C_n}(\mathcal O_X)$, $\operatorname{CM}(\mathcal O_X)$.
\end{enumerate}
\end{remarks}
